\section{Conclusions}
\label{sec:conclusion}

This paper set out to make published hazard detectors usable on the cameras
cities already operate, and to measure honestly what a shared temporal
validation layer contributes toward that goal. The conclusions the
measurements support are stated below in their final, qualified form.

The system stands: five detectors from four independent sources run behind
one declarative gate with enforced viewpoint envelopes, one temporal
validator confirmed to contain no hazard-specific branch, a privacy boundary
with no biometric code path, and an AGPL-free runtime of four packages —
each property held by a build-failing guard, not a promise. The
validator's value, after the controls of \cref{sec:threshold-matched}, is
scoped to what the measurements support: it reaches an oracle-tuned
threshold's false-alarm rate \emph{without} access to the negative
distribution, it deduplicates alerts as no per-frame threshold can, and it
transfers across hazards by configuration — at a closed-form latency cost
(\cref{eq:latency}) and a measured recall cost that thresholding, where it
suffices, does not pay. Two cascade levels that failed their measurements
were removed, not kept (\cref{sec:removed-levels}).
\paragraph{Data and code availability.} Every figure in this paper is
generated from committed measurement files by the build described in
the released repository, including the training and
evaluation scripts that reproduce each number, accompanies this paper.

\paragraph{Acknowledgments.} This work runs on other people's science, named
throughout: the PyroNear association and the
French SDIS; HPWREN at UC San Diego, whose imagery appears here under its
attribution requirement; the U.S. Geological Survey's streamgage camera
network; the authors of ATLANTIS, FloodNet, SeaDronesSee, DRespNeT and
V-FloodNet; and the maintainers of ONNX Runtime, NumPy, OpenCV and
torchvision.
